\section{Introduction}

The United States Medical Licensing Examination (USMLE) represents one of the most consequential assessment systems in medical education worldwide. Administered in three steps, the examination determines physician licensure eligibility and plays a central role in residency program selection across the United States and Canada. Step 1 evaluates the basic biomedical sciences—including anatomy, physiology, biochemistry, pharmacology, pathology, and microbiology—through 280 multiple-choice items administered over a single day. Step 2 Clinical Knowledge (CK) assesses clinical medicine and the application of medical knowledge to patient care scenarios, emphasizing diagnosis, management, and preventive medicine across major clinical disciplines \parencite{usmle2024}.

Pass rates for first-time takers of Step 1 hover around 94\% for US and Canadian medical school graduates, but fall substantially for international medical graduates (IMGs), who comprise approximately 25\% of the active physician workforce in the United States \parencite{nrmp2023}. The disparity reflects not only differences in medical curriculum preparation but also the inherently US-centric framing of many examination questions, which presuppose familiarity with the American healthcare system, insurance structures, referral pathways, and epidemiological patterns that differ markedly from healthcare contexts in other countries \parencite{boulet2006}.

The past three years have witnessed a dramatic proliferation of large language models (LLMs) capable of engaging with complex clinical reasoning tasks. Systems such as OpenAI's GPT-4 \parencite{openai2023gpt4}, Anthropic's Claude \parencite{anthropic2024claude3}, Google's Gemini \parencite{anil2023gemini}, and Meta's Llama \parencite{touvron2023llama2} have each demonstrated strong performance on standardized medical examinations in controlled studies \parencite{nori2023gpt4usmle, kung2023chatgptusmle}. Medical students have rapidly adopted these tools: surveys consistently show that a majority of preclinical and clinical students report using AI chatbots for self-study, question explanation, and concept review \parencite{huh2023medstudents}.

Yet the literature contains critical gaps. First, the majority of published evaluations focus on a single model or compare two models, making it difficult to draw calibrated conclusions about relative strengths across the rapidly expanding ecosystem of frontier LLMs. Second, most studies restrict their evaluation to Step 1 questions or use benchmarks that conflate both steps without reporting disaggregated performance. Given that Step 2 CK has become increasingly important for residency matching—and that many IMG applicants who pass Step 1 struggle disproportionately with the clinical reasoning demands of Step 2 CK—this gap in the literature has practical consequences. Third, and most importantly, no published study has examined LLM performance through the lens of the IMG student, analyzing whether these models exhibit differential accuracy on questions that embed US healthcare system assumptions that would be unfamiliar to graduates of international medical schools.

This paper addresses these gaps through a systematic, reproducible evaluation of eight frontier LLMs on stratified samples of USMLE Step 1 and Step 2 CK questions, with dedicated analysis of the IMG perspective. Specifically, we address four research questions:

\begin{enumerate}
    \item How accurately do frontier LLMs answer USMLE Step 1 and Step 2 CK questions, and how does performance vary by medical subject and difficulty level?
    \item Are there statistically significant performance differences between models after correction for multiple comparisons?
    \item Do models show differential performance on Step 1 versus Step 2 CK questions, and does the pattern of subject-level strengths and weaknesses differ across models?
    \item Do LLMs show lower accuracy on questions containing US-centric clinical assumptions—and does this differential accuracy gap suggest that IMG students who rely on these tools for Step 2 CK preparation may receive systematically biased guidance?
\end{enumerate}

The significance of this work extends beyond academic benchmarking. As IMGs increasingly turn to AI tools for affordable, accessible board preparation—often in contexts where commercial question banks and tutoring services are prohibitively expensive—understanding the reliability and systematic limitations of these tools is essential for informed study practices. A model that performs well on average but fails differentially on questions involving US insurance terminology or American epidemiological baselines could mislead IMG students precisely on the question types where they need the most accurate preparation.

The remainder of this paper is organized as follows. Section 2 reviews related work on LLM performance on medical benchmarks, AI in medical education, and IMG-specific challenges on the USMLE. Section 3 describes our evaluation methodology, including dataset construction, model selection, prompt design, and statistical analysis plan. Section 4 presents our results across overall accuracy, subgroup analyses, and IMG gap analysis. Section 5 discusses the implications of our findings, limitations of the current study, and directions for future research. Section 6 concludes with practical recommendations for medical students and educators.
